\section{Introduction}
\label{sec:introduction}

A judge orders the government to buy a cancer drug within 72 hours. The procurement officer has no time to aggregate demand, no time to research market prices, and faces personal fines if the purchase fails. The drug arrives---but at a 30\% total cost premium over what the same agency pays when no court is watching.

This paper estimates the fiscal costs of such judicial enforcement, exploiting bid-level data on pharmaceutical procurement in S\~ao Paulo, Brazil (2009--2019). Courts and regulatory agencies serve as alternative instruments of social control \citep{glaeser2003rise}, but when courts intervene in procurement, they compress planning timelines and distort incentives in ways that raise prices and reduce competition \citep{coviello2018bidding}. In health care, these costs are especially consequential: judicial mandates directly determine how---and at what price---governments acquire essential medicines \citep{wang2015right, ferraz2009right}. Brazil's 1988 Constitution enshrines health as a ``right of all and a duty of the state,'' and courts have granted virtually every medication request, with 96,000 first-instance claims nationally in 2017 alone---a 913\% increase in S\~ao Paulo since 2008 \citep{cnj2019judicializacao}.

We construct a unique bid-level dataset covering all pharmaceutical procurement by the S\~ao Paulo State Department of Health (SES/SP) from 2009 through 2019, comprising over 479,000 purchase-offer-item observations for medical supplies conducted through the state's electronic procurement platform (BEC).

Our empirical strategy exploits the quasi-experimental variation generated by court orders and administrative requests. These shocks partition purchases into \textit{ordinary} (standard timeline) and \textit{urgent} (compressed timeline). Within urgent purchases, we further distinguish \textit{administrative} requests---which carry no penalty for procurement failure---from \textit{litigated} purchases, where officials face fines, personal liability, and fund seizure. This distinction isolates what we term the ``under the gun'' effect: the additional cost attributable to the threat of judicial punishment. We estimate high-dimensional fixed effects regressions with item, year (or year-month), and public buyer unit (PBU) fixed effects, clustering standard errors at the PBU level.

Urgent purchases carry negotiated prices 5.4\% higher (3.1\% net of quantity), reference prices 2.7\% higher, and attract 5.5\% fewer bidders---yet succeed \textit{more often} (2.1~pp), consistent with officials accepting worse terms to avoid failure. The ``under the gun'' comparison is sharper: litigated tenders cost 23--30\% more than administrative ones that face identical planning constraints but no sanctions. This premium operates entirely through demand fragmentation, not direct price inflation. A placebo on never-litigated items confirms the premium is absent where courts play no role, and supplier fixed effects reveal the cost splits equally between officials accepting worse matches and firms exploiting judicial urgency.

This paper makes three contributions. First, we provide the first estimates of the fiscal cost of judicial enforcement on public procurement outcomes using comprehensive administrative data and high-dimensional fixed effects, contributing to the literature on government spending efficiency \citep{bandiera2009active, best2023procurement, decarolis2018buyer}. Second, we identify and quantify the ``under the gun'' effect by exploiting the institutional distinction between administrative and litigated urgent purchases, contributing to the literature on how accountability mechanisms generate unintended costs \citep{rasul2018management, duflo2018obligations}. Third, we document how judicial mandates distort the planning phase of procurement, adding to the evidence on ex ante design in auction performance \citep{coviello2018bidding, baltrunaite2021discretion}. Heterogeneity analyses (Online Appendix) confirm that enforcement costs are concentrated where competitive pressure is weakest.

The remainder of the paper proceeds as follows. Section~\ref{sec:institutional} describes the institutional setting. Section~\ref{sec:data} presents the data. Section~\ref{sec:empirical} details the empirical strategy. Section~\ref{sec:results} presents results. Section~\ref{sec:conclusion} concludes.
